\documentclass[conference]{IEEEtran}

\usepackage{cite}
\usepackage{amsmath}
\usepackage{booktabs}
\usepackage{array}
\usepackage{url}
\usepackage{hyperref}

\newcommand{\bench}{GPU-NFBench}

\begin{document}

\title{\bench: Benchmarking LLM-Assisted Triage of GPU Numerical Failures in High-Performance ML Software}

\author{
\IEEEauthorblockN{Aryan Shah}
\IEEEauthorblockA{
Texas Academy of Mathematics and Science\\
University of North Texas\\
Denton, TX, USA\\
aryanshah.2431@gmail.com
}
}

\maketitle

\begin{abstract}
High-performance ML software increasingly depends on GPU kernels, graph
compilers, low/mixed-precision arithmetic, and Python-facing accelerator
libraries. When these stacks fail, issue reports about NaNs, overflow, dtype
conversion, precision tolerance, incorrect outputs, compiler failures, and
performance regressions often overlap. This paper presents \bench{}, a
human-labeled benchmark and triage framework for GPU numerical-failure reports.
We construct a 1,191-issue gold benchmark from public GitHub issues across
eight GPU-related repositories using a seven-class taxonomy for invalid values,
overflow/underflow, precision/tolerance, dtype/casting, crash/compile,
performance-only, and non-numerical reports. Weak keyword-derived labels reach
only 50.5\% accuracy, showing that retrieval labels are not sufficient ground
truth. A deterministic voting ensemble reaches 73.8\% accuracy and 0.712 macro
F1; a fine-tuned FLAN-T5-base classifier reaches 80.5\% held-out accuracy and
0.778 macro F1; and an LLM/deterministic agreement mode reaches 89.3\%
accuracy at 68.3\% coverage. Cross-repository transfer remains difficult:
leave-one-repository-out evaluation reaches 66.8\% accuracy. A 119-row blind
audit achieved 80.7\% annotator agreement and Cohen's $\kappa=0.721$, and
follow-on adjudication applied 419 human updates. We release the benchmark,
models, 250-row root-cause evidence-coded extension, error appendix, and
reproducibility artifact with a Zenodo DOI.
\end{abstract}

\begin{IEEEkeywords}
GPU bugs, numerical bugs, benchmark construction, issue mining, LLM-assisted
triage, software reliability, machine learning systems
\end{IEEEkeywords}

\section{Introduction}

Modern AI and scientific workloads rely on GPU software stacks that combine
low-level kernels, graph/compiler transformations, mixed precision, and
high-level Python interfaces. Frameworks such as Triton, PyTorch, JAX, CuPy,
Numba, RAPIDS, and TVM make acceleration accessible, but they also create a
large surface for numerical failures: NaNs, Inf values, silent wrong answers,
overflow, dtype-conversion errors, and compiler/runtime failures that appear
only on specific hardware and software configurations.

Public issue trackers record these failures at scale, but they are noisy. A
search for ``NaN'' can return true invalid-value bugs, missing-data examples,
or unrelated benchmark tables; a search for ``dtype'' can return casting
failures, unsupported features, performance complaints, or compiler crashes.
\bench{} addresses this by separating weak retrieval labels from human gold
labels and by evaluating both full-coverage and abstaining triage systems.

The contributions are:
\begin{itemize}
    \item We construct a 1,191-row human-labeled benchmark of GPU numerical
    failure reports across eight public GPU/HPC software projects.
    \item We define a seven-class taxonomy and preserve normalization and
    adjudication logs, including 419 human updates to the final benchmark.
    \item We evaluate weak labels, lexical baselines, linear classifiers, a
    fine-tuned FLAN-T5-base model, deterministic ensembles, cross-repository
    transfer, chronological transfer, and abstention.
    \item We release the data card, benchmark, model outputs, error appendix,
    50-row human-adjudicated root-cause subset, 250-row evidence-coded
    root-cause extension, and reproducibility artifact.
\end{itemize}

\section{Background and Motivation}

GPU numerical bugs are difficult because symptoms and causes are often
separated. A NaN can originate from overflow in a fused kernel, an unsafe dtype
cast, invalid memory access, compiler lowering, or a tolerance mismatch. Issue
reports rarely state the root cause directly; instead, they contain partial
evidence such as a stack trace, device name, failing kernel, expected CPU
result, tolerance comparison, or maintainer comment. This makes GPU numerical
triage a useful HPEC problem: it combines high-performance GPU software,
low/mixed precision behavior, benchmarking, reliability, and LLM-assisted
analysis.

\section{Benchmark Construction}

The project began with a 930-issue silver seed collected from GPU-related
repositories using terms such as \texttt{nan}, \texttt{inf},
\texttt{overflow}, \texttt{dtype}, \texttt{precision}, \texttt{tolerance},
\texttt{incorrect}, and \texttt{wrong result}. Candidate labels were assigned
from query and text heuristics but were not treated as ground truth. We then
constructed a 191-issue pilot gold benchmark with full public issue/comment
context and expanded it with 1,000 additional human-labeled issue records. The
final v2 benchmark contains 1,191 issues from JAX, PyTorch, RAPIDS cuML,
Triton, RAPIDS cuDF, CuPy, Numba, and Apache TVM. Table~\ref{tab:labels}
shows the primary-label distribution.

\begin{table}[t]
\centering
\caption{Expanded gold benchmark label distribution.}
\label{tab:labels}
\begin{tabular}{lr}
\toprule
Primary label & Issues \\
\midrule
precision\_tolerance & 368 \\
dtype\_casting & 193 \\
crash\_compile & 192 \\
not\_numerical\_failure & 147 \\
nan\_inf & 143 \\
overflow\_underflow & 94 \\
performance\_only & 54 \\
\midrule
Total & 1191 \\
\bottomrule
\end{tabular}
\end{table}

The expanded human annotation file included several labels outside the final
seven-class taxonomy, including API feature requests, compiler/code-generation
descriptions, generic needs-review labels, and performance-regression labels.
Rather than silently discard or recode them, we ran a deterministic
taxonomy-normalization pass and preserved the full change log. In total, 180
rows were mapped into the final taxonomy, and no invalid labels remained. For
example, performance-regression labels were mapped to \texttt{performance\_only};
compiler/code-generation labels were mapped to \texttt{crash\_compile},
\texttt{performance\_only}, or \texttt{not\_numerical\_failure} depending on
the recorded evidence; and API feature requests were mapped to
\texttt{not\_numerical\_failure}. The original labels remain recoverable from
the repair log, so this step is auditable.

We also ran a blind audit sampled from the expansion set. Two annotators
independently labeled 119 rows using the final taxonomy, achieving 80.7\%
agreement and Cohen's $\kappa=0.721$. A third adjudication pass assigned final
labels to those rows, and a second 300-row adjudication pass focused on
low-confidence and model-error examples. Together, these passes applied 419
human updates to the benchmark and changed 211 primary labels. The final
experiments use this canonical v2 benchmark.

\section{Taxonomy}

Each issue receives exactly one primary label:
\begin{itemize}
    \item \textbf{nan\_inf}: invalid-value failures involving NaN, Inf, or
    non-finite outputs.
    \item \textbf{overflow\_underflow}: numerical range failures, including
    saturation, wraparound, or loss to zero.
    \item \textbf{precision\_tolerance}: wrong-answer or tolerance failures
    where approximate arithmetic, accumulation order, or precision is central.
    \item \textbf{dtype\_casting}: dtype promotion, conversion, unsupported
    dtype, or casting semantics failures.
    \item \textbf{crash\_compile}: failures dominated by crashes, compiler
    errors, illegal memory accesses, or build/runtime exceptions.
    \item \textbf{performance\_only}: reports about speed, memory use, or
    scheduling where no numerical correctness issue is present.
    \item \textbf{not\_numerical\_failure}: feature requests, API questions,
    documentation issues, and other non-numerical reports.
\end{itemize}

The taxonomy intentionally separates symptoms from suspected causes. For
example, an issue may be labeled \texttt{nan\_inf} as the primary symptom while
secondary evidence points to dtype conversion, masking, or compiler lowering.
This paper evaluates primary-label triage; secondary-cause prediction is left
as a follow-up task.

\section{Models and Evaluation}

We evaluate models on the 1,191-row expanded gold benchmark. Metrics are
accuracy and macro F1 for the seven-class task. Macro F1 is important because
the dataset is imbalanced: precision/tolerance issues are the largest class,
while performance-only reports are the smallest. We also evaluate a binary
gate that predicts whether an issue is a numerical failure at all. The primary
headline setting is stratified five-fold evaluation. We also report
leave-one-repository-out evaluation, where each repository is held out as the
test project and all other repositories are used for training.

The majority baseline predicts the largest class. The weak-label baseline uses
the original candidate label attached during retrieval. Text baselines include
BM25 nearest-neighbor retrieval, multinomial naive Bayes, TF-IDF logistic
regression, TF-IDF linear SVM, and bigram TF-IDF logistic regression. These
baselines are deliberately simple so that improvements can be attributed to
better supervision and calibrated decision rules rather than opaque modeling
choices.

We fine-tuned a standalone FLAN-T5-base sequence-to-sequence classifier on
benchmark-formatted issue text from the v2 canonical benchmark. The split
contains 945 training rows, 123 validation rows, and 123 held-out test rows. To
reduce class imbalance, the training set was balanced by oversampling minority
classes to 2,058 training examples. The selected v2 checkpoint starts from the
previous expanded-gold FLAN-T5-base checkpoint and is retrained for one epoch
on the v2 labels. It predicts one of the seven primary taxonomy labels directly
from issue text.

The best full-coverage deterministic system is a voting ensemble over the
strongest lexical and linear models. For selective triage, the ensemble can
abstain unless at least $k$ voters agree. This setting matches realistic
engineering use: a triage assistant should confidently label straightforward
issues and flag ambiguous reports for human review.

We also evaluate an LLM-assisted agreement mode. This mode answers only when
the deterministic ensemble and standalone LLM predict the same label; otherwise
it abstains. This is stricter than majority voting and is intended for
high-confidence triage queues where false labels are more costly than deferred
labels.

\section{Results}

\subsection{Full-Coverage Classification}

Table~\ref{tab:fullcoverage} reports full-coverage performance. Weak candidate
labels reach only 50.5\% accuracy, confirming that keyword retrieval labels
are not reliable gold labels. Linear text models perform substantially better
with the expanded human-labeled set. The full-coverage voting ensemble reaches
73.8\% accuracy and 0.712 macro F1.

\begin{table}[t]
\centering
\caption{Expanded gold benchmark full-coverage results.}
\label{tab:fullcoverage}
\begin{tabular}{lcc}
\toprule
Model & Accuracy & Macro F1 \\
\midrule
Majority baseline & 0.309 & 0.067 \\
Candidate weak label & 0.505 & 0.426 \\
BM25 nearest neighbor & 0.571 & 0.534 \\
Naive Bayes & 0.610 & 0.501 \\
TF-IDF logistic regression & 0.715 & 0.688 \\
TF-IDF linear SVM & 0.732 & 0.703 \\
Bigram TF-IDF logistic & 0.724 & 0.694 \\
Voting ensemble & \textbf{0.738} & \textbf{0.712} \\
\bottomrule
\end{tabular}
\end{table}

\subsection{Cross-Repository Generalization}

Table~\ref{tab:loro} reports leave-one-repository-out results. This setting is
harder because model vocabularies, issue templates, and library-specific
concepts shift across projects. The best deterministic models in this setting
reach 66.8\% accuracy, with TF-IDF linear SVM reaching 0.648 macro F1 and the
voting ensemble reaching 0.647 macro F1. The result is substantially above the
weak-label baseline but lower than shuffled evaluation, indicating that
cross-project transfer remains a central challenge for GPU numerical-failure
triage.

\begin{table}[t]
\centering
\caption{Leave-one-repository-out evaluation on expanded gold.}
\label{tab:loro}
\begin{tabular}{lcc}
\toprule
Model & Accuracy & Macro F1 \\
\midrule
Candidate weak label & 0.505 & 0.426 \\
BM25 nearest neighbor & 0.493 & 0.478 \\
Naive Bayes & 0.550 & 0.475 \\
TF-IDF logistic regression & 0.653 & 0.630 \\
TF-IDF linear SVM & \textbf{0.668} & \textbf{0.648} \\
Bigram TF-IDF logistic & 0.663 & 0.633 \\
Voting ensemble & \textbf{0.668} & 0.647 \\
\bottomrule
\end{tabular}
\end{table}

Table~\ref{tab:repo_weakness} summarizes the weakest transfer repositories.
CuPy and Numba are the hardest held-out projects: CuPy mixes dtype/API
compatibility with numerical symptoms, while Numba contains CUDA, LLVM, build,
and runtime language that blurs \texttt{crash\_compile} and
\texttt{not\_numerical\_failure}. This supports cross-repository transfer as a
real benchmark challenge rather than a simple shuffled-text classification
task.

\begin{table}[t]
\centering
\caption{Weakest leave-one-repository-out transfer cases.}
\label{tab:repo_weakness}
\begin{tabular}{lccc}
\toprule
Held-out repo & Issues & Best acc. & Ensemble acc. \\
\midrule
numba/numba & 65 & 0.492 & 0.462 \\
cupy/cupy & 82 & 0.500 & 0.439 \\
rapidsai/cudf & 132 & 0.583 & 0.545 \\
\bottomrule
\end{tabular}
\end{table}

\subsection{Chronological Generalization}

We also evaluate a future-facing chronological split. Issues are sorted by
creation time; the older 80\% form the training set and the newer 20\% form the
test set. The split trains on 952 issues through December 17, 2025 and tests on
239 newer issues. Table~\ref{tab:time} shows that bigram TF-IDF logistic
regression reaches 76.6\% accuracy and 0.738 macro F1, while the voting
ensemble reaches 74.5\% accuracy and 0.722 macro F1. These results suggest
that the benchmark supports triage of later issue reports, but also that the
simpler linear model transfers better than the ensemble under temporal shift.

\begin{table}[t]
\centering
\caption{Chronological 80/20 evaluation on expanded gold.}
\label{tab:time}
\begin{tabular}{lcc}
\toprule
Model & Accuracy & Macro F1 \\
\midrule
Candidate weak label & 0.469 & 0.325 \\
Naive Bayes & 0.598 & 0.475 \\
TF-IDF logistic regression & 0.711 & 0.690 \\
TF-IDF linear SVM & 0.732 & 0.697 \\
Bigram TF-IDF logistic & \textbf{0.766} & \textbf{0.738} \\
Voting ensemble & 0.745 & 0.722 \\
\bottomrule
\end{tabular}
\end{table}

\subsection{Standalone LLM Performance}

Table~\ref{tab:llm} compares the standalone LLM with the best deterministic
ensemble. The v2 standalone FLAN-T5-base model reaches 80.5\% held-out test
accuracy and 0.778 macro F1, outperforming the deterministic ensemble on the
same held-out LLM test split. This changes the practical conclusion from the
earlier expanded-gold version: after adjudication repair, the standalone LLM is
not only a simpler deployable artifact, but also the strongest full-coverage
triage model in the final held-out evaluation. As an additional local
zero-shot comparison, we evaluated Ollama \texttt{llama3.2:3b} on the same 123
held-out rows. It reaches only 31.7\% accuracy and 0.322 macro F1, showing
that task-specific fine-tuning is important. A cloud-backed modern model was
not used for the final comparison because exporting held-out benchmark prompts
to an external service was blocked in the local environment.

\begin{table}[t]
\centering
\caption{Standalone LLM and deterministic comparison on v2 labels.}
\label{tab:llm}
\begin{tabular}{lcc}
\toprule
Model & Accuracy & Macro F1 \\
\midrule
Local Llama 3.2 3B zero-shot & 0.317 & 0.322 \\
TF-IDF linear SVM on LLM test split & 0.691 & 0.669 \\
Voting ensemble on LLM test split & 0.683 & 0.654 \\
FLAN-T5-base standalone LLM & \textbf{0.805} & \textbf{0.778} \\
Binary numerical gate, 5-fold & 0.892 & 0.796 \\
\bottomrule
\end{tabular}
\end{table}

\subsection{Selective Triage}

Table~\ref{tab:abstain} shows the accuracy/coverage tradeoff when the ensemble
abstains unless enough voters agree. With at least three agreeing voters, the
system labels 92.6\% of issues at 76.9\% accuracy. With at least four agreeing
voters, it labels 73.5\% of issues at 84.2\% accuracy. With unanimous or
near-unanimous agreement, it reaches 93.9\% accuracy on 30.1\% of issues. The
LLM/deterministic agreement mode provides a complementary high-confidence
setting: it answers 68.3\% of held-out LLM-test examples at 89.3\% accuracy
and 0.844 macro F1.

\begin{table}[t]
\centering
\caption{Selective prediction with abstention.}
\label{tab:abstain}
\begin{tabular}{lccc}
\toprule
Agreement rule & Coverage & Accuracy & Macro F1 \\
\midrule
Vote $\geq$ 2 & 0.999 & 0.739 & 0.713 \\
Vote $\geq$ 3 & 0.926 & 0.769 & 0.744 \\
Vote $\geq$ 4 & 0.735 & 0.842 & 0.803 \\
Vote $\geq$ 5 & 0.301 & 0.939 & 0.733 \\
LLM/ensemble agreement & 0.683 & 0.893 & 0.844 \\
\bottomrule
\end{tabular}
\end{table}

\subsection{Error Analysis}

The largest error families are semantically plausible boundary cases. The most
common confusion is \texttt{dtype\_casting} predicted as
\texttt{precision\_tolerance} (28 errors), followed by \texttt{nan\_inf}
predicted as \texttt{precision\_tolerance} (22 errors),
\texttt{precision\_tolerance} predicted as \texttt{dtype\_casting} (21 errors),
and \texttt{crash\_compile} predicted as \texttt{precision\_tolerance} (20
errors). These errors reflect the structure of GPU numerical reports: dtype
changes, compiler failures, overflow, and NaN symptoms are often reported
through failed numerical tolerance tests. The next largest errors involve
\texttt{not\_numerical\_failure} issues predicted as crash or dtype failures,
usually because issue text contains API type errors or installation stack
traces even when the report is not a numerical defect. This analysis supports
two follow-up tasks: multi-label secondary-cause prediction and explicit
separation of symptom labels from root-cause labels.

The artifact includes a 12-case error appendix with issue snippets and
boundary explanations. Examples include dtype reports predicted as
precision/tolerance because the failure appears as a numerical mismatch,
NaN/Inf reports embedded in broader tolerance failures, and non-bug support
requests predicted as crash or dtype failures because they contain stack-trace
or type-language cues.

The artifact includes 12 concrete error cases with issue snippets and boundary
explanations. It also contains linked-fix evidence for future root-cause work:
105 pilot issues contain at least one fix or root-cause signal, including 39
explicit pull-request URLs, and a public diff fetch retrieved 488 changed-file
rows. We additionally provide a 50-row human-adjudicated root-cause subset and
a 250-row evidence-coded extension. The extension contains the 50
human-adjudicated rows, 55 linked-fix evidence-coded rows, and 145 issue-text
evidence-coded rows; it is released as root-cause supervision for future work,
not as 250 human-adjudicated labels.

\section{Discussion}

The expanded gold set changes the practical conclusion. On the original
191-issue pilot, full-coverage classification remained near 54\%, and the most
useful mode was abstention. With 1,191 human-labeled examples, full-coverage
triage becomes viable: the deterministic ensemble reaches 73.8\% accuracy, the
standalone LLM reaches 80.5\% held-out accuracy, and chronological evaluation
reaches 76.6\%. The standalone LLM is the simplest deployable full-coverage
artifact, while the deterministic ensemble remains useful as a transparent
baseline and confidence mechanism. For HPEC-style workflows, \bench{} can help
maintainers prioritize numerical defects, help researchers sample failure
clusters for root-cause analysis, and provide a supervised evaluation set for
future debugging agents before patch localization or reproducer generation.

\section{Threats to Validity}

The benchmark relies on human labels and deterministic normalization. Although
all repairs are logged and the audit shows substantial agreement, the full
1,191-row benchmark is not double-adjudicated. The dataset is drawn from public
GitHub issues and may underrepresent private industrial bugs, closed-source
compiler failures, or hardware-specific failures not reported publicly.
Repository sizes are uneven, so cross-project transfer should be interpreted
with care. Finally, many issue reports lack minimal reproducers or confirmed
fixes; the primary task classifies reports, while linked-fix root-cause
prediction remains future work.

\section{Conclusion}

This paper presents \bench{}, a 1,191-row human-labeled benchmark for GPU
numerical failure reports in high-performance ML software. Replacing weak
keyword labels with human gold labels makes reliable triage possible: a
standalone FLAN-T5-base model reaches 80.5\% held-out accuracy, a deterministic
ensemble reaches 73.8\%, chronological evaluation reaches 76.6\%, and
LLM/deterministic agreement reaches 89.3\% accuracy at 68.3\% coverage.
Leave-one-repository-out accuracy is 66.8\%, showing that cross-project
transfer remains difficult. The artifact provides a foundation for GPU
reliability triage, benchmark construction, and root-cause prediction.

\section*{Artifact Availability}

The artifact contains the processed benchmark, expanded gold labels,
taxonomy-normalization report, model-training scripts, fine-tuning files,
evaluation tables, generated reports, qualitative examples, linked-fix evidence,
a data card, completed blind audit files, v2 adjudication files, LLM fine-tuning
files, the 50-row human-adjudicated root-cause subset, the 250-row
evidence-coded root-cause extension, cross-repository weakness analysis, local
LLM baseline outputs, a 12-case error appendix, a reproducibility checklist,
and an external-repository candidate pool for future benchmark expansion. The
data card documents intended use, non-use cases, label provenance, and known
limitations. The artifact is released at
\url{https://github.com/jubs-2431/gpu-nfbench/releases/tag/v1.0-conference}.
The archived Zenodo record is available at
\url{https://zenodo.org/records/20242157} with DOI
\url{https://doi.org/10.5281/zenodo.20242157}.

\begin{thebibliography}{9}

\bibitem{triton2019}
P. Tillet, H. Kung, and D. Cox, ``Triton: An intermediate language and compiler
for tiled neural network computations,'' in \emph{Proceedings of the 3rd ACM
SIGPLAN International Workshop on Machine Learning and Programming Languages},
2019.

\bibitem{cupy2017}
R. Okuta, Y. Unno, D. Nishino, S. Hido, and C. Loomis, ``CuPy: A NumPy-compatible
library for NVIDIA GPU calculations,'' in \emph{Proceedings of Workshop on
Machine Learning Systems at NeurIPS}, 2017.

\bibitem{pytorch2019}
A. Paszke et al., ``PyTorch: An imperative style, high-performance deep learning
library,'' in \emph{Advances in Neural Information Processing Systems}, 2019.

\bibitem{jax2018}
J. Bradbury et al., ``JAX: Composable transformations of Python+NumPy programs,''
2018. [Online]. Available: \url{https://github.com/jax-ml/jax}

\bibitem{numba2015}
S. K. Lam, A. Pitrou, and S. Seibert, ``Numba: A LLVM-based Python JIT
compiler,'' in \emph{Proceedings of the Second Workshop on the LLVM Compiler
Infrastructure in HPC}, 2015.

\bibitem{rapids2024}
RAPIDS Development Team, ``RAPIDS: Collection of GPU accelerated data science
libraries,'' 2024. [Online]. Available: \url{https://rapids.ai/}

\bibitem{tvm2018}
T. Chen et al., ``TVM: An automated end-to-end optimizing compiler for deep
learning,'' in \emph{13th USENIX Symposium on Operating Systems Design and
Implementation}, 2018.

\bibitem{t5}
C. Raffel et al., ``Exploring the limits of transfer learning with a unified
text-to-text transformer,'' \emph{Journal of Machine Learning Research}, vol.
21, no. 140, pp. 1--67, 2020.

\bibitem{flan}
H. W. Chung et al., ``Scaling instruction-finetuned language models,'' \emph{J.
Mach. Learn. Res.}, vol. 25, no. 70, pp. 1--53, 2024.

\end{thebibliography}

\end{document}
